\par\noindent\textbf{(1)} 单程相位为 \(\varphi=2\pi nL/\lambda\)，往返相位 \(\delta=2\varphi\)。镜的振幅反射模平方为 \(R\)，总透射振幅为 \(t^2 e^{i\varphi}\sum_{m=0}^{\infty}(Re^{i\delta})^m=t^2e^{i\varphi}/(1-Re^{i\delta})\)，无损镜满足 \(|t|^2=1-R\)。
\par\noindent\textbf{(2)} 取强度平方得 \(\mathcal T(\delta)=(1-R)^2/[1+R^2-2R\cos\delta]\)。用 \(1+R^2-2R\cos\delta=(1-R)^2+4R\sin^2(\delta/2)\)，得到 \(\mathcal T=1/[1+F\sin^2(\delta/2)]\)，其中系数 \(F=4R/(1-R)^2\)。
\par\noindent\textbf{(3)} 纵模条件 \(\delta=2\pi m\)，相邻频率差 \(\Delta\nu_{\rm FSR}=c/(2nL)\)。在 \(R\to1\) 时半峰条件 \(F\sin^2(\delta/2)=1\)，相位半宽约 \(2/\sqrt F\)，全宽约 \(4/\sqrt F\)，精细度 \(\mathcal F\approx\pi\sqrt R/(1-R)\)。
\par\noindent\textbf{(4)} 若腔内吸收或散射存在，往返振幅系数不再恰为 \(R\)，峰透射率低于1；若镜面相位随频率改变，则 \(\Delta\nu_{\rm FSR}\) 要由总往返相位导数确定。理想 \(c/(2nL)\) 只在介质与反射相位可视为常数时成立。
